\section{Security Goal and Claim Discipline}

Ghostext is designed for a narrow communication setting. Alice wants to send Bob a hidden message through ordinary-looking text. Bob shares Alice's prompt, passphrase, and decoding policy. A passive Censor observes only the transmitted text and does not know the prompt or passphrase. Within that setting, the operational security goal is simple to state: the observed cover text should look like ordinary language-model output rather than like an obviously manipulated payload channel.

\paragraph{Meaning here.} We use this phrase in an implementation-aligned sense rather than in the strongest information-theoretic sense from the formal steganography literature. The design goal is that Ghostext should start from the same next-token distribution that ordinary generation would use under the shared prompt and truncation policy, and that any residual differences should come only from deterministic implementation choices that are explicit in the protocol. In the current implementation, those residual differences are the integer quantization needed for finite arithmetic coding and the retokenization-stability filter needed for deterministic replay. This is weaker than a proof that the full text distribution is indistinguishable from natural text, but it is precise enough to connect code-level behavior to the paper's security wording.

\paragraph{Why the packet layer matters for the security story.} The arithmetic-coding layer does not embed raw plaintext directly. It embeds bytes from an authenticated packet that are intended to be computationally indistinguishable from random under the shared passphrase. This distinction matters because it separates plaintext semantics from the carried bitstream. If plaintext bits were embedded directly, payload structure could interact with token selection in a way that is harder to reason about. By packetizing first, Ghostext aims to make the embedding layer respond to a normalized payload source instead of to message-specific structure. At the same time, packet framing gives Bob an integrity check and clear fail-closed behavior when replay state or transmitted text has changed.

\paragraph{What counts as evidence in this paper.} We separate implemented facts from inferences and from future work. Implemented facts include authenticated packet recovery, deterministic replay under matched state, explicit aborts under mismatch, and bounded retry behavior when local contexts are unsuitable for embedding. Inferences include the systems-level claim that packet normalization is preferable to direct plaintext embedding under the same replay model. Future-work items include detector resistance, prompt-known or multi-message steganalysis, direct KL or total-variation measurement over generated continuations, and cross-hardware determinism. This distinction is deliberate: in a steganography paper, overclaiming detector resistance based only on protocol mechanics would be much easier than supporting it with the right experiments.

\paragraph{Why fail-closed behavior is security-relevant.} Recoverability and stealth are different goals, but they interact. A covert channel that silently emits corrupted plaintext under small replay drift is difficult to reason about operationally, because Alice and Bob may believe communication succeeded when it did not. Ghostext therefore chooses to treat replay mismatch as a hard failure. If the prompt, tokenizer behavior, configuration fingerprint, or packet authentication check does not line up, decoding stops before plaintext release. This does not prove better covertness, but it does make the channel's behavior auditable and keeps the paper's recoverability claims tightly aligned with what the implementation can actually guarantee.

\paragraph{What is hidden and what is not.} The current goal is to hide payload semantics inside fluent text, not to hide every side signal around a communication event. Message length affects packet length and usually affects total generated token count. Prompt families constrain style. Retry behavior can change end-to-end latency. A passive text-only censor in our core model does not necessarily observe all of these signals, but a richer traffic observer could. We therefore do not equate near-distribution preservation at the token-selection layer with complete channel indistinguishability at every systems layer.
